% Part: incompleteness
% Chapter: representability-in-q
% Section: sigma1-completeness

\documentclass[../../../include/open-logic-section]{subfiles}

\begin{document}

\olfileid{inc}{inp}{s1c}
\olsection{Kelengkapan \texorpdfstring{$\Sigma_1$}{Sigma-1}}

Sanajan $\Th{Q}$ lan ekstensi-ekstensine kang konsisten lan bisa
diaksiomatisasi ora jangkep, kita wis weruh yen $\Th{Q}$ mbuktekake
akeh fakta dhasar ngenani numeral. Nyatane, iki bisa dikembangake
luwih adoh. Kanggo mangerteni cakupan perkara kang bisa dibuktekake
ing~$\Th{Q}$, kita ngenalake !!{formula}s $\Delta_0$, $\Sigma_1$,
lan $\Pi_1$. Sacara kasar, !!{formula} $\Sigma_1$ awujud
$\lexists[x][!B(x)]$, ing ngendi $!B$ dibangun mung saka penghubung
proposisional lan kuantor winates. Kita bakal nuduhake yen $!A$ iku
!!{sentence} $\Sigma_1$ kang bener ing $\Struct{N}$, mula
$\Th{Q} \Proves !A$ (\olref{thm:sigma1-completeness}).

\begin{defn}
\ollabel{defn:bd-quant}
\emph{!!{formula} eksistensial winates} awujud
$\lexists[x][(x < t \land !A(x))]$ ing ngendi $t$ iku suku sembarang,
kang lumrahe kita tulis $\bexists{x < t}{!A(x)}$.
%
\emph{!!{formula} universal winates} awujud
$\lforall[x][(x < t \lif !A(x))]$ ing ngendi $t$ iku suku sembarang,
kang lumrahe kita tulis $\bforall{x < t}{!A(x)}$.
\end{defn}

\begin{defn}
\ollabel{defn:delta0-sigma1-pi1-frm}
!!{formula} $!B$ iku $\Delta_0$ yen dibangun saka !!{formula}s
atomik mung nganggo penghubung proposisional lan kuantifikasi
winates.
%
!!{formula} $!A$ iku $\Sigma_1$ yen
$!A \ident \lexists[x][!B(x)]$ ing ngendi $!B$ iku $\Delta_0$.
%
!!{formula} $!A$ iku $\Pi_1$ yen
$!A \ident \lforall[x][!B(x)]$ ing ngendi $!B$ iku $\Delta_0$.
\end{defn}

\begin{lem}
\ollabel{lem:q-proves-clterm-id} Upamana $t$ suku katutup saengga
$\Value{t}{N} = n$. Mula $\Th{Q} \Proves \eq[t][\num n]$.
\end{lem}

\begin{proof}
Kita mbuktekake iki kanthi induksi marang keruwetan~$t$. Kanggo kasus
dhasar, $\Value{\Obj 0}{N} = 0$, lan
$\Th{Q} \Proves \eq[\Obj 0][\num 0]$ amarga
$\num 0 \ident \Obj 0$.
%
Kanggo kasus induksi, tetepna $t_1$ lan $t_2$ minangka suku saengga
$\Value{t_1}{N} = n_1$, $\Value{t_2}{N} = n_2$,
$\Th{Q} \Proves \eq[t_1][\num n_1]$, and
$\Th{Q} \Proves \eq[t_2][\num n_2]$.

Mula $\Value{(t_1')}{N} = n_1 + 1$, lan kita nduweni
$\Th{Q} \Proves \eq[t_1'][{\num n_1}']$ miturut aturan identitas
orde kapisan kang ditrapake marang hipotesis induksi lan !!{formula}
$\eq[\num{n_1}'][\num{n_1}']$,
mula kita nduweni $\Th{Q} \Proves \eq[t_1'][\num{n_1 + 1}]$
miturut definisi numeral.

Kanggo panjumlahan kita nduweni
\[
      \Value{(t_1 + t_2)}{N}
    = \Value{t_1}{N} + \Value{t_2}{N}
    = n_1 + n_2.
\]
Miturut hipotesis induksi lan aturan identitas,
$\Th{Q} \Proves \eq[t_1 + t_2][\num{n_1} + t_2]$, and then
$\Th{Q} \Proves \eq[t_1 + t_2][\num{n_1} + \num{n_2}]$
kanthi ngetrapake aturan identitas kaping pindho.
Miturut \olref[inc][req][bre]{lem:q-proves-add},
$\Th{Q} \Proves \eq[\num{n_1} + \num{n_2}][\num{n_1 + n_2}]$,
mula $\Th{Q} \Proves \eq[t_1 + t_2][\num{n_1 + n_2}]$.

Argumentasi kang padha uga lumaku kanggo~$\times$, nganggo
\olref[inc][req][bre]{lem:q-proves-mult}.
%
Amarga iki nyakup kabeh suku katutup aritmetika, kita nduweni
$\Th{Q} \Proves \eq[t][\num n]$ kanggo saben suku katutup~$t$
saengga $\Value{t}{N} = n$.
\end{proof}

\begin{prob}
Buktekna kanthi rinci bagean \olref{lem:q-proves-clterm-id}
kang ngenani~$\times$.
\end{prob}

\begin{lem}
\ollabel{lem:atomic-completeness}
Upamana $t_1$ lan $t_2$ suku katutup. Mula
\begin{enumerate}
\item Yen $\Value{t_1}{N} = \Value{t_2}{N}$,
    mula $\Th{Q} \Proves \eq[t_1][t_2]$.
\item Yen $\Value{t_1}{N} \neq \Value{t_2}{N}$,
    mula $\Th{Q} \Proves \eq/[t_1][t_2]$.
\item Yen $\Value{t_1}{N} < \Value{t_2}{N}$,
    mula $\Th{Q} \Proves t_1 < t_2$.
\item Yen $\Value{t_2}{N} \leq \Value{t_1}{N}$,
    mula $\Th{Q} \Proves \lnot(t_1 < t_2)$.
\end{enumerate}
\end{lem}

\begin{proof}
Kanggo suku $t_1$ lan $t_2$, kita tetepna
$n = \Value{t_1}{N}$ lan $m = \Value{t_2}{N}$.

Upamana $!A \ident t_1 = t_2$. Miturut
\olref{lem:q-proves-clterm-id},
$\Th{Q} \Proves \eq[t_1][\num n]$ lan
$\Th{Q} \Proves \eq[t_2][\num n]$.
Yen $n = m$, mula $\Th{Q} \Proves \eq[\num n][\num m]$ lan
$\Th{Q} \Proves \eq[t_1][t_2]$ miturut transitivitas identitas.
Yen $n \neq m$, mula $\Th{Q} \Proves \eq/[\num n][\num m]$,
lan miturut transitivitas identitas maneh,
$\Th{Q} \Proves \eq/[t_1][t_2]$.

Saiki tetepna $!A \ident t_1 < t_2$. Kanggo loro kasus iki, kita
nggunakake aksioma~$!Q_8$, kang nyatakake
$x < y \liff \lexists[z][\eq[z' + x][y]]$ kanggo saben $x,y$.

Upamana $\Sat{N}{t_1 < t_2}$. Mula ana $k \in \Nat$ saengga
$n + k + 1 = m$. Miturut \olref{lem:q-proves-clterm-id},
$\Th{Q} \Proves \eq[t_1][\num n]$ lan
$\Th{Q} \Proves \eq[t_2][\num m]$, lan miturut bagean kapisan
lema iki, $\Th{Q} \Proves \eq[\num n + {\num k}'][\num m]$.
Miturut transitivitas identitas, mula
$\Th{Q} \Proves \eq[{\num k}' + t_1][t_2]$,
mula $\Th{Q} \Proves \lexists[z][\eq[z' + t_1][t_2]]$.
Miturut arah saka tengen menyang kiwa saka~$!Q_8$,
$\Th{Q} \Proves t_1 < t_2$.

Kosok baline, upamana $\Sat/{N}{t_1 < t_2}$, yaiku $m \leq n$.
%
Kita makarya ing~$\Th{Q}$ lan nganggep $t_1 < t_2$. Miturut arah
saka kiwa menyang tengen saka~$!Q_8$, ana~$z$ saengga
$\eq[z' + t_1][t_2]$. Amarga
$\Th{Q} \Proves \eq[t_1][\num n]$ lan
$\Th{Q} \Proves \eq[t_2][\num m]$, mula
$\eq[z' + \num n][\num m]$.
%
Kanthi induksi ing njaba teori marang~$m$ nganggo~$!Q_5$, kita entuk
$\eq[z' + \num{n - m}][\Obj 0]$.
Yen $m = n$, mula $\eq/[z'][\Obj 0]$, kang menehi kontradiksi
liwat~$!Q_3$. Yen $m < n$, mula
$\eq[(z' + \num{n - m - 1})'][\Obj 0]$ miturut~$!Q_5$ maneh,
kang menehi kontradiksi liwat~$!Q_3$.
Dadi $\Th{Q} \Proves \lnot(t_1 < t_2)$.
\end{proof}

\begin{lem}
\ollabel{lem:bounded-quant-equiv}
Upamana $!A$ iku !!a{formula}, $t$ suku katutup, lan
$k=\Value{t}{N}$. Mula
\begin{enumerate}
\item $\Th{Q} \Proves \bforall{x<t}{!A(x)}$ iff $\Th{Q} \Proves
    !A(\num 0) \land \dots \land !A(\num{k-1})$.
\item $\Th{Q} \Proves \bexists{x<t}{!A(x)}$ iff $\Th{Q} \Proves
    !A(\num 0) \lor \dots \lor !A(\num{k-1})$.
\end{enumerate}
\end{lem}

\begin{proof}
    Kita mbuktekake kasus kuantor universal winates.
    Yen $\Value{t}{N} = 0$, mula sisih kiwa saka
    ekivalensi bisa dibuktekake ing~$\Th{Q}$, amarga ora ana
    $x<\num 0$ miturut \olref[inc][req][min]{lem:less-zero}.
    Semono uga, kita bisa nganggep disjungsi kosong mung minangka
    $\ltrue$, kang uga bisa dibuktekake ing~$\Th{Q}$.
    %
    Mula upamana $\Value{t}{N} = k+1$ kanggo sawijining
    bilangan asli~$k$. Miturut \olref{lem:q-proves-clterm-id}, kita
    bisa nganggep yen lagi nggarap !!a{formula} awujud
    $\bforall{x<\num{k+1}}{!A(x)}$.
    
    Upamana $\Th{Q} \Proves \bforall{x<\num{k+1}}{!A(x)}$,
    lan tetepna $n \leq k$. Amarga
    $\Th{Q} \Proves \num n < \num{k+1}$ miturut
    \olref{lem:atomic-completeness}, kanthi logika kita entuk
    $\Th{Q} \Proves !A(\num n)$. Kanthi ngetrapake fakta iki $k+1$
    kaping kanggo saben $n \leq k$, kita entuk
    $\Th{Q} \Proves !A(\num 0) \land \dots \land !A(\num k)$
    kaya kang dikarepake.
    
    Kanggo arah liyane, upamana $\Th{Q} \Proves
    !A(\num 0) \land \dots \land !A(\num k)$. Nalika makarya ing
    $\Th{Q}$, upamana $x < \num{k+1}$.
    Miturut \olref[inc][req][min]{lem:less-nsucc}, kita nduweni
    $x = \num 0 \lor \dots \lor x = \num k$, mula kanthi logika
    $!A(x)$ lumaku, lan saka iku pratelan universal
    $\bforall{x<\num{k+1}}{!A(x)}$ uga lumaku.
    
    Bukti ekivalensi kanggo !!{formula}s kanthi kuantor
    eksistensial winates padha carane.
\end{proof}

\begin{prob}
Wenehana bukti rinci kanggo kasus eksistensial ing
\olref{lem:bounded-quant-equiv}.
\end{prob}

\begin{lem}
\ollabel{lem:delta0-completeness}
Yen $!A$ iku !!{sentence} $\Delta_0$ kang bener ing
$\Struct{N}$, mula $\Th{Q} \Proves !A$.
\end{lem}

\begin{proof}
Kita mbuktekake iki kanthi induksi marang keruwetan !!{formula}.
%
Kasus dhasar diwenehi dening \olref{lem:atomic-completeness},
mula kita maju menyang langkah induksi. Supaya prasaja, kita
mbagi kasus negasi dadi sawetara subkasus miturut struktur
!!{formula} kang dinegasi.

\begin{enumerate}
\item Upamana $(!A \land !B)$ bener ing $\Struct{N}$,
mula $!A$ lan $!B$ bener ing~$\Struct{N}$.
Miturut hipotesis induksi, $\Th{Q} \Proves !A$ lan
$\Th{Q} \Proves !B$,
mula $\Th{Q} \Proves (!A \land !B)$ kanthi logika.
%
\item Upamana $\lnot (!A \land !B)$ bener ing $\Struct{N}$,
mula $\lnot !A$ utawa $\lnot !B$ bener ing $\Struct{N}$.
Tanpa ngurangi umum, upamana kang kapisan. Miturut
hipotesis induksi $\Th{Q} \Proves \lnot !A$, mula
$\Th{Q} \Proves \lnot (!A \land !B)$ kanthi logika.
%
\item Upamana $(!A \lor !B)$ bener ing $\Struct{N}$, mula
$!A$ bener ing $\Struct{N}$ utawa $!B$ bener ing
$\Struct{N}$. Tanpa ngurangi umum, upamana kang kapisan.
Miturut hipotesis induksi $\Th{Q} \Proves !A$, mula
$\Th{Q} \Proves (!A \lor !B)$ kanthi logika.
%
\item Upamana $\lnot(!A \lor !B)$ bener ing $\Struct{N}$,
mula $\lnot !A$ lan $\lnot !B$ bener ing $\Struct{N}$.
Mula $\Th{Q} \Proves \lnot !A$ lan
$\Th{Q} \Proves \lnot !B$ miturut hipotesis induksi.
Akibate, $\Th{Q} \Proves \lnot(!A \lor !B)$ kanthi logika.
%
\item Upamana $\bforall{x<t}{!A(x)}$ bener
ing~$\Struct{N}$, ing ngendi $t$ iku suku katutup lan
$k=\Value{t}{N}$. Miturut hipotesis induksi lan logika, yen
$!A(\num n)$ bener ing~$\Struct{N}$ kanggo saben
$n < \Value{t}{N}$, mula $\Th{Q} \Proves
!A(\num 0) \land \dots \land !A(\num{k-1})$.
Miturut \olref{lem:bounded-quant-equiv}, saka kene
$\Th{Q} \Proves \bforall{x<t}{!A(x)}$.
%
\item Kasus kuantor eksistensial winates, nalika
kita duwe !!a{sentence} awujud $\bexists{x < t}{!A(x)}$,
padha karo kasus kuantor universal winates.
%
\item Upamana $\lnot \bforall{x<t}{!A(x)}$ bener
ing~$\Struct{N}$, ing ngendi $t$ iku suku katutup.
!!{sentence} iki ekivalen karo !!{sentence}
$\bexists{x<t}{\lnot !A(x)}$, lan ekivalensine bisa
diturunake ing~$\Th{Q}$, mula kita bisa nggunakake
argumentasi kanggo kuantor eksistensial winates.
%
\item Semono uga, upamana $\lnot \bexists{x<t}!A(x)$
bener ing $\Struct{N}$, ing ngendi $t$ iku suku katutup.
!!{sentence} iki ekivalen ing $\Th{Q}$ karo
$\bforall{x<t}{\lnot!A(x)}$, mula kita bisa nggunakake
argumentasi kanggo kuantor universal winates.
%
\item Pungkasan, upamana $\lnot !A$ bener ing $\Struct{N}$.
Kasus kang isih ana mung nalika $!A$ atomik lan nalika
$\lnot !A \ident \lnot\lnot !B$ kanggo sawijining
!!{sentence} $\Delta_0$ $!B$. Yen $!A$ atomik, mula miturut
\olref{lem:atomic-completeness}, $\Th{Q} \Proves \lnot !A$.
Yen $\lnot !A \ident \lnot\lnot !B$, mula kanthi logika
formula iku bisa dibuktekake ekivalen ing~$\Th{Q}$ karo~$!B$,
kang bener ing~$\Struct{N}$ amarga $\lnot !A$ bener
ing~$\Struct{N}$. Mula miturut hipotesis induksi kita nduweni
$\Th{Q} \Proves \lnot !A$.
\end{enumerate}
\end{proof}

\begin{prob}
Wenehana bukti rinci kanggo kasus eksistensial ing
\olref{lem:delta0-completeness}.
\end{prob}

\begin{thm}
\ollabel{thm:sigma1-completeness}
Yen $!A$ iku !!{sentence} $\Sigma_1$ kang bener
ing~$\Struct{N}$, mula $\Th{Q} \Proves !A$.
\end{thm}

\begin{proof}
Yen $\lexists{x}!A(x)$ iku !!{sentence} $\Sigma_1$
kang bener ing~$\Struct{N}$, mula ana bilangan asli~$n$
lan penetapan variabel~$s$ saengga $s(x) = n$ lan
$\Sat{N}{!A(x)}[s]$. Miturut fakta lumrah ngenani
relasi kepuasan, saka kene kita entuk $\Sat{N}{!A(\num n)}$.
Nanging $!A(\num n)$ iku !!{formula} $\Delta_0$, mula
miturut \olref{lem:delta0-completeness} kita nduweni
$\Th{Q} \Proves !A(\num n)$, lan kanthi logika uga
$\Th{Q} \Proves \lexists[x][!A(x)]$.
\end{proof}

\end{document}
